% Generated by scripts/assemble.py; corrections belong in source/corrections.json.
\examhead{52}{上海财经大学 2026 年数学分析试卷}{53}
% p054-b07; source PDF page 54
\par\Needspace{5\baselineskip}
\blockanchor{p054-b07}
\majorquestion{一}解答题. 每题 15 分, 共 90 分.\par
% p054-b08; source PDF page 54
\blockanchor{p054-b08}
\question{1}判断下述命题是否正确并给出理由.\par
% p054-b09; source PDF page 54
\blockanchor{p054-b09}
(1) 设数列 $\{b_n - a_n\}$ 无界, 数列 $\{a_n\}$ 收敛, 则数列 $\{b_n\}$ 无界.\par
% p054-b10; source PDF page 54
\blockanchor{p054-b10}
(2) 设数列 $\{a_{n}\}$ 满足 $\lim_{n\to \infty}a_n = a$ ，且 $a_{n}$ 是集合 $A$ 的上界，则 $a$ 是集合 $A$ 的上界.\par
% p054-b11; source PDF page 54
\blockanchor{p054-b11}
\question{2}求极限  $\lim_{n\to\infty}\frac{1\times3\times5\times\cdots\times(2n-1)}{2\times4\times6\times\cdots\times(2n)}$ .\par
% p054-b12; source PDF page 54
\blockanchor{p054-b12}
\question{3}计算 $\int_0^{+\infty} \frac{1}{(1 + x^2)(1 + x^{2026})} \mathrm{d}x$.\par
% p054-b13; source PDF page 54
\blockanchor{p054-b13}
\question{4}计算 $\iint_{x^2 + y^2 \leq 1} |3x + 4y| \, \mathrm{d}x \, \mathrm{d}y$.\par
% p054-b14; source PDF page 54
\blockanchor{p054-b14}
\question{5}计算曲线积分 $\int_{C} \frac{y \, \mathrm{d}x - x \, \mathrm{d}y}{x^2 + y^2}$, 其中 $C$ 是 $\frac{x^2}{4} + \frac{y^2}{9} = 1$ 从 (2,0) 到 (-2,0) 的一段曲线, 取逆时针方向.\par
% p054-b15; source PDF page 54
\blockanchor{p054-b15}
\question{6}求平面 $x + y + z = 0$ 与椭球面 $x^{2} + y^{2} + 4z^{2} = 1$ 相交而成的椭圆面积.\par
% p054-b16; source PDF page 54
\par\Needspace{5\baselineskip}
\blockanchor{p054-b16}
\majorquestion{二}证明题. 每题 10 分, 共 60 分.\par
% p054-b17; source PDF page 54
\par\Needspace{6\baselineskip}
\blockanchor{p054-b17}
\question{7}设\par
\nopagebreak[4]
% p054-b18; source PDF page 54
\blockanchor{p054-b18}
\displayformula{f (x) = \left\{ \begin{array}{l l} {p,} & {x =  \frac {p}{q}   (p \in \mathbb {N} ^ {+},   q \in \mathbb {Z} \backslash \{0 \},   p, q   \text {互素}),} \\ {1,} & {x = 0,} \\ {0,} & {x   \text {是无理数}.} \end{array} \right.}
% p054-b19; source PDF page 54
\blockanchor{p054-b19}
证明： $f(x)$ 在任意区间上无界.\par
% p054-b20; source PDF page 54
\blockanchor{p054-b20}
\question{8}设 $f(x)$ 在 $[a, b]$ 上有定义, 证明: $f(x)$ 在 $[a, b]$ 上一致连续的充要条件是 $f(x)$ 将 $[a, b]$ 上的任意 Cauchy 数列映射为 Cauchy 数列.\par
% p054-b21; source PDF page 54
\blockanchor{p054-b21}
\question{9}已知 $f(x)$ 在 $[a, b]$ 上二阶可导，且有 $f(a) = f(b) = 0$，若存在 $c \in (a, b)$ 使得 $f(c) > 0$. 证明：存在 $\xi \in (a, b)$，使得 $f''(\xi) < 0$.\par
% p054-b22; source PDF page 54
\par\Needspace{6\baselineskip}
\blockanchor{p054-b22}
\question{10}设 $f(x)$ 在 $[a, +\infty)$ 上三阶可导, 证明: 若 $\lim_{x \to +\infty} f(x), \lim_{x \to +\infty} f'''(x)$ 都存在且有限, 则\par
\nopagebreak[4]
% p054-b23; source PDF page 54
\blockanchor{p054-b23}
\displayformula{\lim _ {x \rightarrow + \infty} f ^ {\prime} (x) = \lim _ {x \rightarrow + \infty} f ^ {\prime \prime} (x) = \lim _ {x \rightarrow + \infty} f ^ {\prime \prime \prime} (x) = 0.}
% p055-b00; source PDF page 55
\par\Needspace{6\baselineskip}
\blockanchor{p055-b00}
\question{11}设\par
\nopagebreak[4]
% p055-b01; source PDF page 55
\blockanchor{p055-b01}
\displayformula{f (x, y) = \left\{ \begin{array}{l l} x y \cos \frac {1}{\sqrt {x ^ {2} + y ^ {2}}}, & (x, y) \neq (0, 0), \\ 0, & (x, y) = (0, 0). \end{array} \right.}
% p055-b02; source PDF page 55
\blockanchor{p055-b02}
证明： $f(x,y)$ 在 $(0,0)$ 处可微，但偏导函数不连续.\par
% p055-b03; source PDF page 55
\blockanchor{p055-b03}
\question{12}设 $a_{n} \geq 0$, 且幂级数 $\sum_{n=1}^{\infty} a_{n} x^{n}$ 的收敛半径为 $R$, 证明: $\lim_{x \to R^{-}} \sum_{n=1}^{\infty} a_{n} x^{n} = \sum_{n=1}^{\infty} a_{n} R^{n}$.\par
